
%===================================
%===================================
\section{Conclusiones}
\label{sec:conclusiones}

\begin{enumerate}
    \item Los resultados muestran que la versión MPI sin optimizaciones del compilador obtiene mejoras bastante limitadas frente a la versión secuencial optimizada. El speedup promedio apenas alcanza valores de 0.87 con 2 procesos, 1.3 con 3 procesos y 1.6 con 4 procesos, lo que evidencia que gran parte del tiempo ganado por paralelismo termina siendo consumido por la sobrecarga asociada a la comunicación, sincronización y distribución de datos entre procesos. En este escenario, la paralelización por sí sola no es suficiente para convertir la aplicación en una solución realmente eficiente con respecto a la optimización por compulador usada en el algoritmo secuencial.

    \item Cuando se habilitan optimizaciones del compilador, el comportamiento cambia de manera significativa. Los speedups promedio aumentan hasta 10.96 con 2 procesos, 14.14 con 3 procesos y 15.52 con 4 procesos, alcanzando incluso un máximo de 29.41x para el caso de \(N = 6400\) con 4 procesos. Esto demuestra que las optimizaciones relacionadas con localidad de memoria, mejor uso de caché, vectorización y reorganización de bucles tienen un impacto decisivo en algoritmos intensivos en memoria como la multiplicación de matrices densas.

    \item Se observa además que el rendimiento mejora a medida que aumenta el tamaño del problema. Para matrices pequeñas, los costos de creación de procesos, sincronización y transferencia de información representan una parte importante del tiempo total de ejecución, generando speedups modestos e incluso inferiores a 1 en algunos casos. Sin embargo, conforme el tamaño de la matriz crece, cada proceso realiza una mayor cantidad de trabajo computacional y el costo fijo de coordinación se amortiza, permitiendo que el paralelismo sea mucho más efectivo en tamaños como \(N = 3200\) y \(N = 6400\).

    \item Aunque el rendimiento aumenta al utilizar más procesos, la escalabilidad no es perfectamente lineal. Pasar de 2 a 4 procesos mejora el tiempo total de ejecución, pero el rendimiento no se duplica proporcionalmente. Esto sugiere la existencia de limitaciones asociadas al acceso a memoria, la comunicación MPI y el balance de carga entre procesos. Además, el tiempo total medido incluye operaciones adicionales como lectura desde NFS, distribución mediante \texttt{MPI\_Scatterv}, recolección con \texttt{MPI\_Gatherv} y escritura final de resultados, factores que también afectan la eficiencia global.

    \item Desde el punto de vista computacional, los resultados son coherentes con la complejidad \(O(n^3)\) del algoritmo de multiplicación de matrices. A medida que \(n\) aumenta, el costo del cálculo crece mucho más rápido que la sobrecarga de coordinación entre procesos, favoreciendo así el aprovechamiento del paralelismo. Las optimizaciones del compilador reducen considerablemente el costo del kernel principal, permitiendo que cada proceso ejecute una mayor cantidad de trabajo útil y logrando un mejor aprovechamiento del hardware disponible.

    \item Finalmente, la versión sin optimizaciones presenta una variabilidad considerablemente mayor en los tiempos de ejecución, especialmente en tamaños pequeños y en configuraciones con pocos procesos. Esto indica que los costos fijos del sistema y el overhead de MPI tienen un peso importante cuando la carga computacional todavía es reducida. En contraste, la versión optimizada muestra un comportamiento mucho más estable y consistente, reflejando una utilización más eficiente de los recursos de cómputo y memoria.
    \item Una de las conclusiones más importantes es que aumentar la cantidad de procesos no garantiza automáticamente una mejora en el rendimiento. Los resultados muestran que, cuando el núcleo de multiplicación no está optimizado, gran parte del tiempo se pierde en coordinación entre procesos, sincronización y acceso a memoria. En otras palabras, para que MPI realmente aporte beneficios, primero es necesario que cada proceso tenga suficiente trabajo útil y que el código base sea eficiente.

    \item Los experimentos también evidencian que el principal cuello de botella del problema no está en la capacidad de cálculo del procesador, sino en el acceso a memoria y el manejo de caché. La gran diferencia de rendimiento entre la versión optimizada y la no optimizada demuestra que optimizaciones como el mejor orden de los bucles y el aprovechamiento de la localidad de memoria tienen un impacto mucho mayor de lo esperado. Esto confirma que, en multiplicación de matrices densas, un código bien optimizado puede ser más importante que simplemente agregar más procesos.

    \item En términos de escalabilidad, la implementación presenta un comportamiento favorable conforme aumenta el tamaño de la matriz. Para tamaños pequeños, el overhead de MPI y las operaciones de coordinación representan una parte importante del tiempo total, por lo que las mejoras son limitadas. Sin embargo, en matrices grandes como \(N = 3200\) y \(N = 6400\), el costo del cálculo domina sobre la comunicación y el sistema logra aprovechar mejor el paralelismo disponible, especialmente en la versión optimizada.

\end{enumerate}